
\documentclass[10pt,aspectratio=169]{beamer}
\usepackage[utf8]{inputenc}
\usepackage[T1]{fontenc}
\usepackage[catalan]{babel}
\usepackage{amsmath,amssymb,booktabs,array,multicol}
\usepackage{tikz,pgfplots,xcolor,listings}
\usetikzlibrary{positioning,arrows.meta,calc,shapes.geometric}
\pgfplotsset{compat=1.18}
\usepgfplotslibrary{fillbetween}

\definecolor{UABGreen}{RGB}{0,103,71}
\definecolor{UABLight}{RGB}{224,241,236}
\definecolor{DeepBlue}{RGB}{20,65,110}
\definecolor{AccentOrange}{RGB}{230,126,34}
\definecolor{SoftGray}{RGB}{245,247,250}
\definecolor{CodeBack}{RGB}{248,248,248}
\definecolor{CodeFrame}{RGB}{210,215,220}

\usetheme{Madrid}
\usecolortheme[named=UABGreen]{structure}
\setbeamercolor{title}{fg=white,bg=UABGreen}
\setbeamercolor{frametitle}{fg=white,bg=UABGreen}
\setbeamercolor{block title}{fg=white,bg=DeepBlue}
\setbeamercolor{block body}{fg=black,bg=SoftGray}
\setbeamercolor{alerted text}{fg=AccentOrange}
\setbeamertemplate{navigation symbols}{}
\setbeamertemplate{itemize item}{\color{UABGreen}$\blacktriangleright$}
\setbeamertemplate{itemize subitem}{\color{AccentOrange}$\bullet$}
\setbeamertemplate{enumerate item}{\color{UABGreen}\insertenumlabel.}
\setbeamertemplate{footline}{%
  \leavevmode%
  \hbox{%
  \begin{beamercolorbox}[wd=.72\paperwidth,ht=2.4ex,dp=1ex,leftskip=1.4em]{author in head/foot}%
    Estadística -- Grau en Enginyeria Informàtica
  \end{beamercolorbox}%
  \begin{beamercolorbox}[wd=.28\paperwidth,ht=2.4ex,dp=1ex,center]{date in head/foot}%
    \insertframenumber{} / \inserttotalframenumber
  \end{beamercolorbox}}%
  \vskip0pt%
}

\lstdefinelanguage{R}{%
  morekeywords={if,else,repeat,while,function,for,in,next,break,TRUE,FALSE,NULL,NA,NaN,Inf,
    library,mean,median,sd,var,summary,table,prop.table,xtabs,addmargins,round,
    ggplot,aes,geom_point,geom_line,geom_col,geom_histogram,geom_density,geom_smooth,
    geom_tile,facet_wrap,labs,theme_minimal,summarise,group_by,mutate,read_csv,filter,
    arrange,set.seed,rbinom,rgeom,rpois,runif,rexp,rnorm,dnorm,pnorm,qnorm,dt,pt,qt,
    dchisq,pchisq,qchisq,df,pf,qf,dbinom,pbinom,dpois,ppois,replicate,sample,rep,seq,
    tibble,data.frame,confint,t.test,var.test,chisq.test,binom.test,prop.test,lm,
    broom,glance,augment,tidy,future_map_dbl,plan,multisession,map_dbl,purrr,future,furrr},
  sensitive=true,
  morecomment=[l]\#,
  morestring=[b]",
  morestring=[b]'
}
\lstset{language=R,basicstyle=\ttfamily\scriptsize,keywordstyle=\color{DeepBlue}\bfseries,
  commentstyle=\color{gray},stringstyle=\color{AccentOrange},backgroundcolor=\color{CodeBack},
  frame=single,rulecolor=\color{CodeFrame},breaklines=true,columns=fullflexible,keepspaces=true,
  showstringspaces=false,
  literate={à}{{\`a}}1 {è}{{\`e}}1 {é}{{\'e}}1 {í}{{\'i}}1 {ï}{{\"i}}1 {ò}{{\`o}}1 {ó}{{\'o}}1 {ú}{{\'u}}1 {ü}{{\"u}}1 {ç}{{\c{c}}}1 {À}{{\`A}}1 {È}{{\`E}}1 {É}{{\'E}}1 {Í}{{\'I}}1 {Ò}{{\`O}}1 {Ó}{{\'O}}1 {Ú}{{\'U}}1}

\newcommand{\R}{\textsf{R}}
\newcommand{\RStudio}{\textsf{RStudio/Posit}}
\newcommand{\GitHub}{\textsf{GitHub}}
\newcommand{\important}[1]{\textcolor{UABGreen}{\textbf{#1}}}
\newcommand{\exemple}[1]{\textcolor{AccentOrange}{\textbf{#1}}}
\newcommand{\Prob}{\mathbb{P}}
\newcommand{\E}{\mathbb{E}}
\newcommand{\Var}{\operatorname{Var}}
\newcommand{\IC}{\operatorname{IC}}

\title[Inferència III]{Inferència III}
\author{Estadística -- Grau en Enginyeria Informàtica}
\date{Curs 2026--27}

\begin{document}
\begin{frame}[plain]
  \titlepage
  \vspace{-0.45cm}
  \begin{center}
  \small\textcolor{DeepBlue}{Estadística computacional amb \R, simulació i exemples d'enginyeria informàtica}
  \end{center}
\end{frame}


\begin{frame}[fragile]{De l'estimació a la decisió}
\begin{itemize}
  \item Els intervals de confiança responen: ``quin rang és plausible?''
  \item Els tests d'hipòtesi responen: ``les dades són compatibles amb una afirmació concreta?''
\end{itemize}
\begin{block}{Exemple informàtic}
Una versió nova redueix la latència mitjana per sota de 100 ms? Una taxa d'error supera el 1\%?
\end{block}
\end{frame}

\begin{frame}[fragile]{Objectius}
\begin{enumerate}
  \item Formular hipòtesi nul·la i alternativa.
  \item Entendre error de tipus I, tipus II i significació.
  \item Calcular estadístics de contrast i p-valors.
  \item Distingir tests bilaterals i unilaterals.
  \item Aplicar tests simples per mitjana, proporció i variància amb \R.
\end{enumerate}
\end{frame}

\begin{frame}[fragile]{Hipòtesi nul·la i alternativa}
\begin{columns}[T]
\column{0.50\textwidth}
\[
H_0: \theta=\theta_0
\]
\[
H_1: \theta\ne\theta_0,\quad \theta>\theta_0,\quad \theta<\theta_0.
\]
\begin{itemize}
  \item $H_0$ és el model de referència.
  \item $H_1$ representa l'efecte que busquem evidenciar.
\end{itemize}
\column{0.44\textwidth}
\begin{block}{Exemples}
\begin{itemize}
  \item $H_0: \mu=100$ ms.
  \item $H_1: \mu<100$ ms.
  \item $H_0: p=0.01$.
  \item $H_1: p>0.01$.
\end{itemize}
\end{block}
\end{columns}
\end{frame}

\begin{frame}[fragile]{Errors i p-valor}
\begin{itemize}
  \item Error tipus I: rebutjar $H_0$ quan és certa.
  \item Error tipus II: no rebutjar $H_0$ quan $H_1$ és certa.
  \item Significació $\alpha$: risc màxim tolerat d'error tipus I.
\end{itemize}
\begin{block}{p-valor}
Probabilitat d'obtenir un resultat igual o més extrem que l'observat, assumint que $H_0$ és certa.
\end{block}
\begin{alertblock}{Regla}
Si p-valor $<\alpha$, rebutgem $H_0$.
\end{alertblock}
\end{frame}

\begin{frame}[fragile]{Cues del test}
\begin{columns}[T]
\column{0.33\textwidth}
\begin{block}{Cua dreta}
$H_1:\theta>\theta_0$\\
Ex.: la taxa d'error és massa alta.
\end{block}
\column{0.33\textwidth}
\begin{block}{Cua esquerra}
$H_1:\theta<\theta_0$\\
Ex.: la latència ha baixat.
\end{block}
\column{0.33\textwidth}
\begin{block}{Bilateral}
$H_1:\theta\ne\theta_0$\\
Ex.: el rendiment ha canviat.
\end{block}
\end{columns}
\vspace{0.3cm}
\begin{alertblock}{Clau}
La hipòtesi alternativa es decideix abans de mirar el resultat.
\end{alertblock}
\end{frame}

\begin{frame}[fragile]{Test simple per a la mitjana}
\begin{block}{$\sigma$ coneguda o $n$ gran}
\[
E=\frac{\bar X-\mu_0}{\sigma/\sqrt n}\sim N(0,1)\quad \text{sota }H_0.
\]
\end{block}
\begin{lstlisting}
xbar <- 96; mu0 <- 100; sigma <- 15; n <- 64
Eobs <- (xbar - mu0)/(sigma/sqrt(n))
pval_esq <- pnorm(Eobs)  # H1: mu < mu0
Eobs; pval_esq
\end{lstlisting}
\end{frame}

\begin{frame}[fragile]{Test simple per a la mitjana amb $\sigma$ desconeguda}
\begin{block}{t-test}
\[
T=\frac{\bar X-\mu_0}{S/\sqrt n}\sim t_{n-1}.
\]
\end{block}
\begin{lstlisting}
set.seed(1)
lat <- rnorm(25, mean = 96, sd = 15)
t.test(lat, mu = 100, alternative = "less")
\end{lstlisting}
\begin{itemize}
  \item \texttt{alternative = "less"}: $H_1:\mu<100$.
  \item \texttt{alternative = "greater"}: $H_1:\mu>100$.
  \item Per defecte, \texttt{two.sided}.
\end{itemize}
\end{frame}

\begin{frame}[fragile]{Test per a una proporció}
\begin{columns}[T]
\column{0.52\textwidth}
Suposem que el sistema ha de tenir una taxa d'error màxima de 1\%.
En 5000 peticions observem 68 errors.
\[
H_0:p=0.01,\qquad H_1:p>0.01.
\]
\column{0.42\textwidth}
\begin{lstlisting}
x <- 68; n <- 5000
prop.test(x, n, p = 0.01,
          alternative = "greater",
          correct = FALSE)
binom.test(x, n, p = 0.01,
           alternative = "greater")
\end{lstlisting}
\end{columns}
\end{frame}

\begin{frame}[fragile]{Test per a una variància}
\begin{itemize}
  \item Volem saber si la variabilitat supera un llindar.
  \item Si la població és normal:
\end{itemize}
\[
\chi^2=\frac{(n-1)S^2}{\sigma_0^2}\sim\chi^2_{n-1}\quad \text{sota }H_0.
\]
\begin{lstlisting}
n <- 30; s <- 18; sigma0 <- 15
Eobs <- (n-1)*s^2/sigma0^2
pval <- 1 - pchisq(Eobs, df = n-1) # H1: sigma^2 > sigma0^2
Eobs; pval
\end{lstlisting}
\end{frame}

\begin{frame}[fragile]{Resultat estadístic vs decisió pràctica}
\begin{block}{No n'hi ha prou amb el p-valor}
Cal mirar també:
\begin{itemize}
  \item mida de l'efecte,
  \item interval de confiança,
  \item cost d'actuar o no actuar,
  \item qualitat i independència de les dades.
\end{itemize}
\end{block}
\begin{alertblock}{En enginyeria}
Una diferència estadísticament significativa pot ser irrellevant si només redueix 0.1 ms, i una diferència no significativa pot requerir més dades.
\end{alertblock}
\end{frame}

\end{document}
